\subsection{Differential operators}

14.1.1. Let \(k\) be an integer such that \(0 \leqslant k \leqslant r\) , and let \(\mathbf{P}^k(\mathbf{E})\) be the bundle of jets of sections of order \(k\) from E (12.6.1). Set

\[
\mathbf {D} ^ {k} (\mathbf {E}, \mathbf {F}) = \mathscr {L} (\mathbf {P} ^ {k} (\mathbf {E}); \mathbf {F}).
\]

It is a vector bundle with base X and class \(C^{r-k}\) . A section of this bundle over an open set U of X is identified with a map

\[
\mathbf {D}: \mathrm{P} ^ {k} (\mathrm{E}) | \mathrm{U} \rightarrow \mathrm{F} | \mathrm{U}
\]

which commutes with the projection onto U and is linear on each fiber; such a section is called a differential operator on U, of type E \(\to\) F, and of order \(\leqslant k\) . Let \(h \in N_{K} \cup \{0\}\) , with \(0 \leqslant h \leqslant r-k\) . We denote by \(\mathcal{D}_{\mathrm{U}}^{k,h}(\mathrm{E}, \mathrm{F})\) the set of differential operators on U, of type E \(\to\) F and of order \(\leqslant k\) , which are of class \(C^{h}\) (as sections of \(\mathrm{D}^{k}(\mathrm{E}, \mathrm{F})\) , or as morphisms of \(\mathrm{P}^{k}(\mathrm{E})|\mathrm{U}\) into \(\mathrm{F}|\mathrm{U}\), it amounts to the same thing). If \(\mathrm{M} = \mathrm{D}^{k}(\mathrm{E}, \mathrm{F})\) , one has \(\mathcal{D}_{\mathrm{U}}^{k,h}(\mathrm{E}, \mathrm{F}) = \mathcal{S}_{\mathrm{M}}^{h}(\mathrm{U})\) ; it is a module over \(\mathcal{C}^{h}(\mathrm{U})\) , cf. 7.4.1.

If \(0 \leqslant k' \leqslant k\) , the morphism \(r^{k, k'}: \mathrm{P}^{k}(\mathrm{E}) \to \mathrm{P}^{k'}(\mathrm{E})\) (cf. 12.6.6) defines an injection from \(\mathrm{D}^{k'}(\mathrm{E}, \mathrm{F})\) to \(\mathrm{D}^{k}(\mathrm{E}, \mathrm{F})\) ; every differential operator of order \(\leqslant k'\) is thus identified with a differential operator of order \(\leqslant k\) . A differential operator of order \(\leqslant k\) which is not of order \(\leqslant k - 1\) is sometimes said to be of order \(k\) . If \(h \in \mathbf{N}_{\mathbb{K}} \cup \{0\}\) and \(h \leqslant r - k\) , and if U is an open subset of X, we have:

\[
0 \subset \mathcal {D} _ {\mathrm{U}} ^ {0, h} (\mathrm{E}, \mathrm{F}) \subset \mathcal {D} _ {\mathrm{U}} ^ {1, h} (\mathrm{E}, \mathrm{F}) \subset \dots \subset \mathcal {D} _ {\mathrm{U}} ^ {k, h} (\mathrm{E}, \mathrm{F}).
\]

If \(r=\infty\) or \(\omega\) , one denotes by \(\mathcal{D}_{\mathrm{U}}^{\infty,h}(\mathrm{E},\mathrm{F})\) the union of the \(\mathcal{D}_{\mathrm{U}}^{k,h}(\mathrm{E},\mathrm{F})\) for \(k\geqslant0\) ; an element of this union is called a differential operator on U of bounded order (of type E \(\rightarrow\) F, of class C \(^{h}\) ) or sometimes simply a differential operator on U.

14.1.2 (“Operators of order zero”). We have \(\mathbf{P}^0 (\mathbf{E}) = \mathbf{E}\) (12.5.1) and \(\mathrm{D}^0 (\mathrm{E},\mathrm{F}) = \mathcal{L}(\mathrm{E};\mathrm{F})\) . If U is open in X, and if \(h\in \mathbf{N}_{\mathbb{K}}\cup \{0\}\) and \(h\leqslant r\) , an element of \(\mathcal{D}_{\mathbf{U}}^{0,h}(\mathbf{E},\mathbf{F})\) is a morphism of class \(\mathbf{C}^h\) from E\textbar U to F\textbar U.

14.1.3 (“Weakening of structure”). Suppose K = R, and let \(r' \in N_{R}\) with \(r' \leqslant r\) . Let X', E', and F' be the manifold and the vector bundles of class \(C^{r'}\) obtained by weakening the structure from X, E, and F respectively. Let k be an integer such that \(0 \leqslant k \leqslant r'\) . Then \(P^{k}(E')\) is a bundle of class \(C^{r'-k}\) deduced from \(P^{k}(E)\) by weakening the structure, and one has an analogous result for \(D^{k}(E', F')\) . In particular, if U is open in X, and if \(0 \leqslant h \leqslant r - k'\) , one has \(\mathcal{D}_{\mathrm{U}}^{k,h}(E', F') = \mathcal{D}_{\mathrm{U}}^{k,h}(E, F)\) .

14.1.4. Let D be a differential operator on X, of type E \(\to\) F, of order \(\leqslant k\) (where \(0 \leqslant k \leqslant r\)). Let U be an open subset of X, and let s be a section of class \(\mathbf{C}^{m}\) of E over U, with \(m \in N _{K} \cup \{0\}\) and \(k \leqslant m \leqslant r\). Then \(j^{k}(s)\) is a section of class \(\mathbf{C}^{m-k}\) of \(P^{k}(E)\) over U (12.6.1); its image under D is denoted \(D_{U}(s)\), or D(s). It is a section of F over U. Suppose D is of class \(\mathbf{C}^{h}\) , where \(h=m-k\). Then \(D_{U}(s)\) is of class \(\mathbf{C}^{h}\) and the maps

\[
\mathrm{D} _ {\mathrm{U}} \colon \mathcal {S} _ {\mathrm{E}} ^ {m} (\mathrm{U}) \rightarrow \mathcal {S} _ {\mathrm{F}} ^ {h} (\mathrm{U})
\]

thus obtained enjoy the following properties:

\begin{enumerate}
\def\labelenumi{(\arabic{enumi})}
\item
  \(D_{U}\) is K-linear.
\item
  For every \(s \in \mathcal{S}_{\mathbb{E}}^{m}(\mathrm{U})\) and every open set \(\mathrm{V}\) of \(\mathrm{U}\) , one has \(\mathrm{D}_{\mathrm{V}}(s|\mathrm{V}) = \mathrm{D}_{\mathrm{U}}(s)|\mathrm{V}\) .
\item
  For every \(s \in \mathcal{S}_{\mathbb{E}}^{m}(\mathrm{U})\) and every \(x \in \mathrm{U}\) such that \(j_{x}^{k}(s) = 0\) , one has \(D_{\mathrm{U}}(s)(x) = 0\) .
\end{enumerate}

(3') For every coordinate system \(\xi = (\xi^1, \ldots, \xi^n)\) in U and any frame \(s = (s_1, \ldots, s_d)\) of E over U (7.4.4), there exist sections \(n_{i,\alpha}(1 \leqslant i \leqslant d, \alpha \in \mathbf{N}^n, |\alpha| \leqslant k)\) of F on U, of class \(C^h\) and such that, for every family \((f_i)_{1 \leqslant i \leqslant d}\) of elements of \(\mathcal{C}^m(U)\) , one has

\[
\mathrm{D} _ {\mathrm{U}} \left(\sum_ {1 \leqslant i \leqslant d} f _ {i}. s _ {i}\right) = \sum_ {1 \leqslant i \leqslant d, | \alpha | \leqslant k} \Delta_ {\xi} ^ {\alpha} (f _ {i}). n _ {i, \alpha},
\]

\begin{enumerate}
\def\labelenumi{(\arabic{enumi})}
\setcounter{enumi}{3}
\tightlist
\item
  For any function \(f \in \mathcal{C}^m(\mathrm{U})\) and any mapping \(\theta\) of \(\mathcal{S}_{\mathrm{U}}^{m}(\mathrm{E})\) to \(\mathcal{S}_{\mathrm{U}}^{h}(\mathrm{F})\) , denote by \(\operatorname{ad}(f)\theta\) the map
\end{enumerate}

\[
s \mapsto f. \theta (s) - \theta (f. s)
\]

of \(\mathcal{S}_{\mathrm{U}}^{m}(\mathrm{E})\) to \(\mathcal{S}_{\mathrm{U}}^{h}(\mathrm{F})\) . Then, for every function \(f \in \mathcal{C}^{m}(\mathrm{U})\) , there exists a differential operator \(\mathbf{L}\) on \(\mathbf{U}\) , of type \(\mathrm{E} \to \mathrm{F}\) of order \(\leqslant k - 1\) and of class \(\mathbf{C}^{h}\) such that

\[
(\operatorname{ad} (f) \mathrm{D} _ {\mathrm{U}}) (s) = \mathrm{L} _{\mathrm{U}}(s) \quad \text { for all } s \in \mathscr{S}_{\mathrm{E}}^{m}(\mathrm{U}).
\]

\begin{enumerate}
\def\labelenumi{(\arabic{enumi})}
\setcounter{enumi}{4}
\tightlist
\item
  One has
\end{enumerate}

\[
\operatorname{ad} (f _ {0}) \dots \operatorname{ad} (f _ {k}) \mathrm{D} _ {\mathrm{U}} = 0 \quad \text { for all } f _ {0}, \dots , f _ {k} \text { in } \mathscr {C} ^ {m} (\mathrm{U}).
\]

If we set I = \{0, . . ., k\} and \(f_{H} = \prod_{i \in H} f_i\) for every subset H of I, the preceding relation is equivalent to:

\[
\sum_ {\mathrm{H} \in \mathrm{I}} (- 1) ^ {\operatorname{Card} (\mathrm{H})} f _ {\mathrm{H}}. \mathrm{D} _ {\mathrm{U}} (f _ {\mathrm{I-H}}. s) = 0
\]

for all \(f_{0},\ldots,f_{k}\) in \(\mathcal{C}^{m}(\mathrm{U})\) and \(s\) in \(\mathcal{S}_{\mathbf{E}}^{m}(\mathrm{U})\) .

14.1.5 (“Characterization of differential operators”). Let \(k, m\) and \(h = m - k\) be as in 14.1.4. For every open set U of X, let \(\mathbf{D}_{\mathbf{U}}\) be a map from \(\mathcal{S}_{\mathbf{E}}^{m}(\mathbf{U})\) to \(\mathcal{S}_{\mathbf{F}}^{h}(\mathbf{U})\) . Suppose that conditions 1, 2, 3 (resp. 1, 2, 3') of 14.1.4 are satisfied for every open subset U of X. Then there exists a unique element D of \(\mathcal{D}_{\mathbf{X}}^{k,h}(\mathbf{E},\mathbf{F})\) such that the \(\mathbf{D}_{\mathbf{U}}\) are the corresponding maps. \(^{(1)}\) In what follows, we identify D with the family of the \(\mathbf{D}_{\mathbf{U}}\) .

When \(m = \infty\) or \(\omega\) , one has \(h = m\) and we may replace conditions (3) and (3') by one of the conditions (4) and (5).

14.1.6 (“Scalar operators”). Suppose that E and F are equal to the trivial bundle \(K_{X}\) . A differential operator of type \(E \rightarrow F\) is then called a scalar differential operator or simply a differential operator; if it is of order \(\leqslant k\) , it is a section of the bundle \(\mathcal{L}(\mathrm{P}^{k}(\mathrm{X}); \mathrm{K}_{\mathrm{X}}) = \mathrm{P}^{k}(\mathrm{X})^{*}\) , dual of the bundle \(\mathrm{P}^{k}(\mathrm{X})\) ; using the isomorphism

\[
i ^ {- 1} \colon \mathrm{P} ^ {k} (\mathrm{X}) ^ {*} \rightarrow \mathrm{T} ^ {(k)} (\mathrm{X}) \quad (\text { cf. } 13.2.5),
\]

we see that a scalar differential operator of order \(\leqslant k\) is identified with a section of the vector bundle \(\mathbf{T}^{(k)}(\mathbf{X})\) , i.e. with a field of point distributions of order \(\leqslant k\) .

Take in particular for X an open set of \(K^{n}\) , where n is an integer \(\geqslant0\) . The fields of point distributions

\[
\Delta^ {\alpha} \colon x \mapsto \Delta_ {x} ^ {\alpha} \quad (\text { cf. } 1 3. 2. 6)
\]

are scalar differential operators of class \(\mathbf{C}^{\omega}\) . For every \(h\in\mathbf{N}_{\mathbf{K}}\) , the \(\Delta^{\alpha}\) for \((|\alpha|\leqslant k)\) form a basis of the \(\mathcal{C}^{h}(\mathbf{X})\)-module \(\mathcal{D}_{\mathbf{X}}^{k,h}(\mathbf{K}_{\mathbf{X}},\mathbf{K}_{\mathbf{X}})\) .

14.1.7 (“Complex differential operators on a real manifold”). Suppose K = R and that the vector bundles E and F are equipped with complex structures (8.8.1); let k be an integer such that \(0 \leqslant k \leqslant r\) . The bundle \(\mathrm{P}^{k}(\mathrm{E})\) is then equipped with a complex structure (12.6.3); the bundle \(\mathcal{L}_{\mathbf{C}}(\mathrm{P}^{k}(\mathrm{E}); \mathrm{F})\) (7.8.5) of complex linear maps from \(\mathrm{P}^{k}(\mathrm{E})\) to F is a vector sub-bundle of \(\mathcal{L}(\mathrm{P}^{k}(\mathrm{E}); \mathrm{F}) = \mathrm{D}^{k}(\mathrm{E}, \mathrm{F})\) ; we denote it \(\mathrm{D}_{\mathbf{C}}^{k}(\mathrm{E}, \mathrm{F})\) . A section of \(\mathrm{D}_{\mathbf{C}}^{k}(\mathrm{E}, \mathrm{F})\) is called a complex differential operator, of type \(E \to F\) and of order \(\leqslant k\) . If \(h \in N_{K}\) and \(h \leqslant r - k\) , and if U is an open set of X, we denote similarly by \(\mathcal{D}_{\mathbf{U}}^{k,h}(\mathrm{E}, \mathrm{F})_{\mathbf{C}}\) the space of sections of class \(C^{h}\) of \(\mathrm{D}_{\mathbf{C}}^{k}(\mathrm{E}, \mathrm{F})\) on U; it is a C-vector space. The other definitions and results of this paragraph extend analogously to complex differential operators; we leave their formulation to the reader.

14.1.8 (“Composition”). Let G be a vector bundle of class \(C^{r}\) with base X. Let \(k'\) and \(k''\) be positive integers with sum \(k \leqslant r\) , and let \(D'\) (resp. \(D''\) ) be a differential operator on X, of type \(E \to F\) (resp. of type \(F \to G\) ), of order \(\leqslant k'\) (resp. \(\leqslant k''\) ). Suppose \(D': P^{k'}(E) \to F\) is of class \(C^{h'}\) , with \(h' \in N_{K} \cup \{0\}\) and \(k'' \leqslant h' \leqslant r - k'\) ; the map

\[
\mathbf {D} _ {*} ^ {\prime} = \mathbf {P} ^ {k ^ {\prime \prime}} (\mathbf {D} ^ {\prime}) \colon \mathbf {P} ^ {k ^ {\prime \prime}} (\mathbf {P} ^ {k ^ {\prime}} (\mathbf {E})) \rightarrow \mathbf {P} ^ {k ^ {\prime \prime}} (\mathbf {F})
\]

This map is of class \(\mathbf{C}^{h'-k''}\) (12.6.3). Let \(\beta\) be the canonical homomorphism of \(\mathbf{P}^{k}(\mathbf{E})\) to \(\mathbf{P}^{k''}(\mathbf{P}^{k'}(\mathbf{E}))\) , cf. 12.6.5. By composing the maps

\[
\mathrm{P} ^ {k} (\mathrm{E}) \xrightarrow {\beta} \mathrm{P} ^ {k ^ {\prime \prime}} (\mathrm{P} ^ {k ^ {\prime}} (\mathrm{E})) \xrightarrow {\mathrm{D} _ {*} ^ {\prime}} \mathrm{P} ^ {k ^ {\prime \prime}} (\mathrm{F}) \xrightarrow {\mathrm{D} ^ {\prime \prime}} \mathrm{G},
\]

we obtain a differential operator of type E → G and of order ≤ k. This operator is called the composite of D'' and D'; it is denoted by D'' ∘ D'.

\(^{1}\) It suffices, moreover, that condition (3′) be satisfied for a family \((\xi_{\lambda})\) of coordinate systems and a family of frames \((s_{\lambda})\) whose domains cover X.

If U is an open subset of X, and if \(s \in \mathcal{S}_{\mathrm{E}}^{k}(\mathrm{U})\) , one has

\[
(\mathbf {D} ^ {\prime \prime} \circ \mathbf {D} ^ {\prime}) (s) = \mathbf {D} ^ {\prime \prime} (\mathbf {D} ^ {\prime} (s))
\]

(where both sides are sections of G over U); this justifies the terminology and notation adopted.

Suppose now \(\mathbf{D}^{\prime \prime}\) is of class \(\mathbf{C}^{h^{\prime \prime}}\) , with \(h^{\prime \prime} \in \mathrm{N}_{\mathbb{K}} \cup \{0\}\) and \(h^{\prime \prime} \leqslant r-k^{\prime\prime}\) . Then \(\mathbf{D}^{\prime \prime} \circ \mathbf{D}^{\prime}\) is of class \(\mathbf{C}^{h}\) , with \(h=\inf(h'-k'',h'')\) .

Suppose we have E = F = G, and that \(k' \leqslant h''\) , in which case \(D' \circ D''\) is defined. Then \(D' \circ D'' - D'' \circ D'\) is a differential operator of type \(E \to E\) and of order \(\leqslant k - 1\) ; we denote it \([D', D'']\) .

14.1.9 ("Associativity"). With the notations and hypotheses as in 14.1.8, let H be a vector bundle of class \(\mathbf{C}^r\) and base X, and let \(\mathbf{D}''\) be a differential operator on X, of type \(\mathbf{G} \to \mathbf{H}\) and of order \(\leqslant k^{m}\) , with \(k' + k'' + k'' \leqslant r\) . We assume that \(\mathbf{D}'\) is of class \(\mathbf{C}^{h'}\) , with \(h' \in \mathrm{N}_{\kappa} \cup \{0\}\) and \(k'' + k'' \leqslant h' \leqslant r - k'\) , and that \(\mathbf{D}''\) is of class \(\mathbf{C}^{h''}\) , with \(h'' \in \mathrm{N}_{\kappa} \cup \{0\}\) and \(k'' \leqslant h'' \leqslant r - k''\) . Then the composites

\[
\mathbf {D} ^ {\prime \prime \prime} \circ (\mathbf {D} ^ {\prime \prime} \circ \mathbf {D} ^ {\prime}) \quad \text { and } \quad (\mathbf {D} ^ {\prime \prime \prime} \circ \mathbf {D} ^ {\prime \prime}) \circ \mathbf {D} ^ {\prime}
\]

are defined and are equal.

14.1.10. Examples

\begin{enumerate}
\def\labelenumi{\alph{enumi})}
\tightlist
\item
  Let U be an open subset of X, let \(D \in \mathcal{D}_{\mathrm{U}}^{k,h}(\mathrm{E}, \mathrm{F})\) with \(h \in N_{K} \cup \{0\}\) and \(h \leqslant r-k\) , and let \(f \in \mathcal{C}^{h+k}(\mathrm{U})\) (resp. \(f \in \mathcal{C}^{h}(\mathrm{U})\) ). Multiplication by f in E (resp. F) is a differential operator on U of order \(\leqslant0\) , and of type E \(\to\) E (resp. F \(\to\) F), cf. 14.1.2. The composites \(D\circ f\) and \(f\circ D\) are defined and belong to \(\mathcal{D}_{\mathrm{U}}^{k,h}(\mathrm{E},\mathrm{F})\) ; we set
\end{enumerate}

\[
\operatorname{ad} (f) \mathrm{D} = f \circ \mathrm{D} - \mathrm{D} \circ f \in \mathcal {D} _ {\mathrm{U}} ^ {k - 1, h} (\mathrm{E}, \mathrm{F}),
\]

cf. 14.1.4, (4). One thus endows \(\mathcal{D}_{\mathrm{U}}^{k,h}(\mathrm{E},\mathrm{F})\) with the structure of a \(\mathcal{C}^{h+k}(\mathrm{U})\)-module on the right (resp. of a \(\mathcal{C}^h (\mathrm{U})\)-module on the left). The left module structure coincides with that of 14.1.1.

\begin{enumerate}
\def\labelenumi{\alph{enumi})}
\setcounter{enumi}{1}
\item
  Suppose \(r = \infty\) or \(r = \omega\) . The differential operators of type \(E \to E\) and of class \(C^r\) form an associative K-algebra with identity element.
\item
  Take for X an open subset of \(K^{n}\) . The scalar differential operators \(\Delta^{\alpha}\) (14.1.6) satisfy the formulas
\end{enumerate}

\[
\Delta^ {\alpha} \circ \Delta^ {\beta} = ((\alpha , \beta)) \Delta^ {\alpha + \beta}.
\]

They commute with each other. Denote by \(D_{i}\) the operator \(\Delta^{\varepsilon_{i}}\) ("i-th partial derivative"). If K has characteristic 0, we have

\[
\Delta^ {\alpha} = \frac {1}{\alpha !} D _ {1} ^ {\alpha_{1}} \dots D _ {n} ^ {\alpha_{n}}.
\]

If K has characteristic \(p \neq 0\) , one has \(\mathbf{D}_{i}^{p} = 0\) for all \(i\) .

14.1.11 ("Multidifferential operators"). Let \(E_{1}, \ldots, E_{n}\) be vector bundles of class \(C^{r}\) with base X, and let \(k, k_{1}, \ldots, k_{n}\) be integers belonging to \([0, r]\) . Set

\[
\mathrm{L} (k _ {1}, \dots , k _ {n}) = \mathscr {L} (\mathrm{P} ^ {k _ {1}} (\mathrm{E} _ {1}) \otimes \dots \otimes \mathrm{P} ^ {k _ {n}} (\mathrm{E} _ {n}); \mathrm{F}).
\]

If \(k_i \leqslant k\) for all \(i\) , the projections \(r^k, k_i: \mathrm{P}^k(\mathrm{E}_i) \to \mathrm{P}^{k_i}(\mathrm{E}_i)\) make it possible to identify the bundle \(\mathrm{L}(k_1, \ldots, k_n)\) with a subbundle of the bundle \(\mathrm{L}(k) = \mathrm{L}(k, \ldots, k)\) . There exists a smallest subbundle \(\mathrm{M}(k)\) of \(\mathrm{L}(k)\) containing all the subbundles \(\mathrm{L}(k_1, \ldots, k_n)\) , for \(k_1 + \cdots + k_n = k\) . A section \(\mathrm{H}\) of \(\mathrm{M}(k)\) is called an n-differential operator (or multidifferential operator) of type \((\mathrm{E}_1, \ldots, \mathrm{E}_n) \to \mathrm{F}\) and of order \(\leqslant k\) . If \(s_i\) ( \(1 \leqslant i \leqslant n\) ) is a section of class \(\mathbf{C}^m (m \in \mathbb{N}_{\mathbb{K}}, m \geqslant k)\) of \(\mathbf{E}_i\) on an open set U of X, one defines \(\mathrm{H}_{\mathrm{U}}(s_1, \ldots, s_n) = \mathrm{H}(s_1, \ldots, s_n)\) as the image under H of the section \(j^k(s_1) \otimes \cdots \otimes j^k(s_n)\) of the bundle

\[
\mathbf {P} ^ {k} (\mathrm{E} _ {1}) \otimes \dots \otimes \mathbf {P} ^ {k} (\mathrm{E} _ {n}).
\]

It is a section of F\textbar U. If H is of class \(\mathbf{C}^h\) , this section is of class \(\mathbf{C}^p\) with \(p = \inf(m - k, h)\) .

When \(E_{1}, \ldots, E_{n}\) and F are equal to \(K_{x}\) , one says that H is a scalar n-differential operator (or multidifferential operator).

Example. --- Let J be a finite set, X an open subset of K \(^{J}\) . Let H be a scalar n-differential operator of order \(\leqslant k\) on X. There exists a unique family of scalar functions

\[
c (\alpha (1), \dots , \alpha (n)) _ {(\alpha (1), \dots , \alpha (n)) \in \mathbf {N} ^ {\mathrm{J}}}, \quad \sum_ {i = 1} ^ {n} | \alpha (i) | \leqslant k
\]

on X such that

\[
\mathrm{H} _ {\mathrm{X}} \left(f _ {1}, \dots , f _ {n}\right) = \sum_ {\alpha (1), \dots , \alpha (n)} c (\alpha (1), \dots , \alpha (n)) \Delta^ {\alpha (1)} \left(f _ {1}\right) \dots \Delta^ {\alpha (n)} \left(f _ {n}\right)
\]

for every family \((f_1, \ldots, f_n)\) of functions of class \(\mathbf{C}^k\) on \(X\) . For \(H\) to be of class \(\mathbf{C}^h\) , it is necessary and sufficient that the functions \(c(\alpha(1), \ldots, \alpha(n))\) be of class \(\mathbf{C}^h\) .

